\documentclass[11pt,a4paper,oneside]{scrartcl}
\usepackage[utf8]{inputenc}
\usepackage[english,russian]{babel}
\usepackage[top=1cm,bottom=1cm,left=1cm,right=1cm]{geometry}

\begin{document}
\pagestyle{empty}

\begin{center}
{\large\scshape\bfseries Программа курса <<Математическая логика>>}\\
{\large\scshape Вопросы ко второму коллоквиуму.}\\
\itshape ИТМО, группы M3232--M3239, осень 2026 г.
\end{center}

%\vspace{0.3cm}

\begin{enumerate}
\item Исчисление предикатов:
\begin{enumerate}
\item Категорические силлогизмы. Термины, предикат, субъект, фигуры, 
модусы (сильные, слабые, неправильные),
ограничения, контрпримеры на ограничения.
\item Язык исчисления предикатов. Функциональные и предикатные символы. Кванторы. 
Метапеременные, правила расстановки скобок, сокращения записи.
\item Вхождения, свободные вхождения, подстановка, свобода для подстановки.
\item Предметное множество, оценки, функции и предикаты.
\item Общезначимость, следование.
\item Теория доказательств. Правила вывода. Выводимость.
\item Теорема о дедукции в исчислении предикатов. Ограничения.
Контрпример, демонстрирующий важность ограничений.
\item Лемма о перестановке подстановки и оценки. Теорема о корректности исчисления предикатов.
Контрпример, демонстрирующий важность ограничений.
\end{enumerate}
\item Теорема Гёделя о полноте исчисления предикатов.
\begin{enumerate}
\item Непротиворечивые множества формул (замкнтуые, с кванторами и бескванторные).
\item Пополнение множества формул. Полные непротиворечивые множества формул.
\item Существование моделей у непротиворечивых множеств формул в 
бескванторном исчислении предикатов. Примеры.
\item Поверхностные кванторы (предварённая форма). Эквивалентность формул 
формулам с поверхностными кванторами.
\item Сколемизация.
\item Теорема Гёделя о полноте исчисления предикатов (формулировка).
\item Полнота исчисления предикатов (идея доказательства).
\item Теорема Гёделя о компактности.
\end{enumerate}
\item Неразрешимость исчисления предикатов
\begin{enumerate}
\item Машина Тьюринга: определения. Тезис Чёрча.
\item Язык. Разрешимость языка. Разрешимость теории. 
\item Задача об останове, её неразрешимость. 
\item Неразрешимость исчисления предикатов (идея доказательства).
\end{enumerate}
\item Представление (целых, радицональных, вещественных) чисел через натуральные.
Дедекиндовы сечения.
\item Арифметизация математики.
\begin{enumerate}
\item Подход Лейбница к формализации категорических силлогизмов.
\item Аксиоматика Пеано. 
\item Арифметические операции (сложение, умножение) в аксиоматике Пеано.
\item Порядок теории (0, 1, 2). Теории первого порядка. Примеры теорий второго порядка.
\item Формальная арифметика. 
\item Примитивно-рекурсивные и рекурсивные функции. Функция Аккермана. 
\item Функции вычисления простых чисел. Частичный логарифм.
\item Выразимость отношений и представимость функций в формальной арифметике. Характеристические функции.
\item Бета-функция Гёделя. 
Гёделева нумерация. Рекурсивность представимых в формальной арифметике функций.
\end{enumerate}
\item Теоремы Гёделя о неполноте Ф.А.
\begin{enumerate}
\item Непротиворечивость (эквивалентные определения), $\omega$-не\-про\-ти\-во\-ре\-чи\-вость. 
\item Первая теорема Гёделя о неполноте арифметики. Формулировка первой теоремы Гёделя о 
неполноте арифметики в форме Россера. 
\item Синтаксическая и семантическая неполнота арифметики.
\item Неполнота расширений формальной арифметики.
\item Ослабленные варианты: арифметика Пресбургера, система Робинсона.
\item Вторая теорема Гёделя о неполноте арифметики, $Consis$. 
\item Лемма об автоссылках. Условия Гильберта-Бернайса-Лёба. 
\item Неразрешимость формальной арифметики. Теорема Тарского о невыразимости истины.
\end{enumerate}
\end{enumerate}
\end{document}
